\section{The determinant in the full representation space}
\label{sec:c3-collision-determinant}

A collision of two point modules changes both the critical coefficient
complex and the choices of a constituent line. A nontrivial punctual
extension has one invariant line with multiplicity two. A scalar
two-dimensional module has a projective line of invariant lines.
The potential distinguishes these two collisions even though their
underlying point supports agree.

Work over $\mathbb C$ with rational constructible complexes on analytic
quotient stacks. Let
\[
 V=\operatorname{Mat}_2(\mathbb C)^3,\qquad
 f(A,B,C)=\operatorname{Tr}(A[B,C]),\qquad
 \mathfrak X=[V/GL_2].
\]
The action is simultaneous conjugation. There is no additional torus
scaling the spatial coordinates. Let $\mathfrak M$ be the classical
critical substack of $f$. Its coefficient complex is
\[
 \mathcal P=\left.\phi_f^p\mathbb Q_{\mathfrak X}[8]
 \right|_{\mathfrak M},\qquad \phi_f^p=\phi_f[-1].
\]
Use the ambient critical-chart orientation. The dimension of
$\mathfrak X$ is $12-4=8$. A geometric stalk means pullback along a
chosen object $\operatorname{Spec}\mathbb C\to\mathfrak M$.
It differs from equivariant cohomology over the object's stabilizer.

Write $A=aI+A_0$, $B=bI+B_0$, and $C=cI+C_0$, where the subscript
denotes the traceless part. Fix the basis
\[
 H=\begin{pmatrix}1&0\\0&-1\end{pmatrix},\qquad
 E=\begin{pmatrix}0&1\\0&0\end{pmatrix},\qquad
 F=\begin{pmatrix}0&0\\1&0\end{pmatrix}
\]
of $\mathfrak{sl}_2$. Let $M\in\operatorname{Mat}_3(\mathbb C)$ have
the coordinate columns of $A_0,B_0,C_0$ in this basis.

Two one-parameter families separate the collisions. The family
\[
 A_0(s)=\begin{pmatrix}0&1\\s&0\end{pmatrix},\qquad B_0=C_0=0,
\]
has two eigenlines for $s\ne0$ and a nonzero nilpotent matrix at
$s=0$. The family $A_0(t)=tE$, $B_0=C_0=0$, then approaches the
scalar triple at $t=0$. Both families lie in the full critical
scheme, and neither requires diagonalizing a repeated eigenvalue.

\begin{proposition}[The full determinant model]
\label{prop:c3-full-determinant}
The potential and its classical critical scheme are
\begin{equation}
 f=2\det M,\qquad
 \operatorname{Crit}(f)=\mathbb C^3\times Y,
 \quad Y=\{M:\operatorname{rank}M\leq1\}.
 \label{eq:c3-full-determinant}
\end{equation}
The action of $GL_2$ is trivial on $\mathbb C^3$ and acts on the
rows of $M$ through its adjoint representation, of determinant one.
The locus $Y^*=Y\setminus\{0\}$ is smooth of dimension five.
Its only singular boundary is the origin.
\end{proposition}

\begin{proof}
Scalar matrices disappear from the commutator and from its trace.
The form $\operatorname{Tr}(A_0[B_0,C_0])$ is alternating in the
three arguments, by cyclicity of trace. Since $[E,F]=H$ and
$\operatorname{Tr}(H^2)=2$, it is twice the determinant in the chosen
basis. Its nine partial derivatives are twice the signed $2\times2$
minors. Those minors generate the reduced rank-at-most-one
determinantal scheme: the parametrization $M=u v^t$ identifies its
coordinate ring with the subring generated by $u_i v_j$ in
$\mathbb C[u_1,u_2,u_3,v_1,v_2,v_3]$.
The relations are the $2\times2$ minors, as follows by exchanging
indices in monomials with the same row and column multiplicities.
This also proves reducedness.

On $Y^*$, the vectors $u,v$ are nonzero and are unique up to
$(u,v)\mapsto(tu,t^{-1}v)$. Thus $Y^*$ is the smooth quotient
of $(\mathbb C^3\setminus0)^2$ by this free $\mathbb G_m$ action.
Its dimension is $3+3-1=5$. At the origin all nine tangent directions
survive, since the defining minors are quadratic. The origin is
therefore singular. Conjugation preserves the alternating trace form,
so its action on $\mathfrak{sl}_2$ has determinant one.
\end{proof}

The rank-one expression can be written
\[
 (A_0,B_0,C_0)=(v_xN,v_yN,v_zN),\qquad N\ne0,
 \quad v\ne0.
\]
If $-\det N\ne0$, the matrix $N$ is semisimple. Its two eigenlines
give two distinct point supports. The first coordinates can agree
when $v_x=0$, while another coordinate still separates the points.
If $\det N=0$, then $N^2=0$ and $N$ has rank one. The module is a
nontrivial extension of a point module by itself. Its common invariant
line is $\ker N=\operatorname{im}N$.
The remaining stratum $M=0$ consists of scalar triples, or direct
sums of two copies of the same point module.

The stabilizers are $\mathbb G_m^2$ in the semisimple case,
$\{sI+tN:s\ne0\}\cong\mathbb G_m\times\mathbb G_a$ in the
nilpotent case, and $GL_2$ at a scalar triple. These statements follow
by solving $gN=Ng$. In particular, a scalar collision increases the
stabilizer dimension from two to four.

\section{Transverse and scalar Milnor fibres}
\label{sec:c3-collision-stalks}

\begin{lemma}[Quadratic normal form away from the scalar stratum]
\label{lem:c3-collision-morse}
At every point of $Y^*$, the full determinant has a nondegenerate
quadratic normal form in four variables. Its Milnor fibre has reduced
rational cohomology $\mathbb Q$ in degree three and zero in every
other degree. The normal vanishing-cycle local system on $Y^*$ is
trivial.
\end{lemma}

\begin{proof}
Choose a nonzero entry as pivot and permute rows and columns to
place it in the upper left corner. Up to the resulting constant
sign, write
\[
 M=\begin{pmatrix}t&r\\s&B\end{pmatrix},\qquad
 Z=B-t^{-1}sr,
 \quad t\ne0.
\]
Block elimination gives $\det M=t\det Z$. The entries of $t,r,s$
are five independent coordinates along $Y^*$, and its equations in
this chart are $Z=0$. Multiplying the first row of $Z$ by the unit
$2t$, including the permutation sign, gives
\[
 f=w_{11}w_{22}-w_{12}w_{21}.
\]
This change of coordinates is invertible on the pivot chart.
It fixes the function exactly and requires no square root of $t$.
The displayed Hessian has rank four. Its nonzero fibre is isomorphic
to $SL_2(\mathbb C)$, which retracts by polar decomposition to
$SU_2\cong S^3$. This proves the stalk calculation.

The local rank-one vanishing-cycle spaces form a local system.
The principal $\mathbb C^*$ bundle
$(\mathbb C^3\setminus0)^2\to Y^*$ has simply connected total
space and connected fibre. The homotopy exact sequence gives
$\pi_1(Y^*)=0$. Consequently this rational local system is trivial.
Its generator can be fixed on one pivot chart by the complex Thom
orientation of an isotropic plane in the quadratic normal form.
The orientation then extends over $Y^*$.
\end{proof}

\begin{theorem}[The scalar collision stalk]
\label{thm:c3-scalar-milnor}
At a scalar triple, the Milnor fibre of the full potential has the
homotopy type of $SU_3$. Its rational cohomology ring is
\begin{equation}
 H^*(F_0,\mathbb Q)=\bigwedge_{\mathbb Q}(\alpha_3,\alpha_5),
 \qquad |\alpha_3|=3,\quad |\alpha_5|=5.
 \label{eq:c3-collision-milnor-cohomology}
\end{equation}
The geometric stalks of the stack-normalized complex $\mathcal P$
are therefore
\[
\begin{array}{c|c|c}
\text{stratum}&\text{nonzero stalk degrees}&\text{dimensions}\\\hline
M\in Y^*&-4&1\\
M=0&-4,-2,1&1,1,1.
\end{array}
\]
Milnor monodromy acts as the identity on each of these cohomology
groups. These are stalk statements and do not assert a splitting
of the global equivariant complex.
\end{theorem}

\begin{proof}
Translate the three scalar coordinates to zero. Contracting them
does not change $f$ and decreases the norm. It remains to compute
the local fibre of $2\det M$ at the origin of $\mathbb C^9$.
For $\epsilon\ne0$, scalar multiplication by a chosen cube root
of $\epsilon/2$ identifies this fibre in a ball with
\[
 SL_3(\mathbb C)\cap\{\|M\|<R\},
 \qquad R=\delta/|\epsilon/2|^{1/3},
\]
using the Frobenius norm on the nine coordinates and ball radius
$\delta$. For sufficiently small $|\epsilon|$, one has $R>\sqrt3$.

Write $M=UP$ by polar decomposition, with $U\in SU_3$ and $P$
positive Hermitian of determinant one. The homotopy $UP\mapsto UP^s$,
$1\geq s\geq0$, stays in $SL_3$. It also stays in the ball.
Indeed, if the eigenvalues of $P$ are $e^{x_i}$, then
$\sum_i x_i=0$ and
\[
 \|UP^s\|^2=\sum_{i=1}^3 e^{2s x_i}.
\]
This function is convex, and its derivative at $s=0$ is zero.
It increases for $s\geq0$, so decreasing $s$ decreases the norm.
The homotopy retracts the local fibre onto $SU_3$.
The same determinant-fibre description is recorded in
\cite[\S1, pp.~4--5]{DamonDeterminantFibres}.

The first-column map $SU_3\to S^5$ has fibre $SU_2\cong S^3$.
In its rational Serre spectral sequence, only bidegrees
$(0,0),(0,3),(5,0),(5,3)$ occur. Every differential has zero source
or zero target. Thus the cohomology has rank one in degrees
$0,3,5,8$. Multiplicativity identifies the degree-eight class with
the product of the degree-three and degree-five generators.
Their squares vanish by graded commutativity in characteristic zero.
This proves the cohomology ring formula.

For an unshifted constant complex, vanishing-cycle stalk cohomology
is the reduced cohomology of the Milnor fibre.
The convention $\phi^p=\phi[-1]$ and the stack shift $[8]$ give
\[
 H^j(\mathcal P)_x=\widetilde H^{j+7}(F_x,\mathbb Q).
\]
This yields the displayed degrees. Equivalently, on the smooth
atlas $V\to\mathfrak X$ of relative dimension four,
the perverse complex $\phi_f^p\mathbb Q_V[12]$ has degrees
$-8,-6,-3$ at a scalar triple. Pullback of $\mathcal P$ is that
complex shifted by $[-4]$.

At the origin a geometric Milnor monodromy is $M\mapsto\zeta M$,
where $\zeta=e^{2\pi i/3}$. Left multiplication by the path
\[
 \operatorname{diag}(e^{2\pi it/3},e^{2\pi it/3},e^{-4\pi it/3}),
 \qquad 0\leq t\leq1,
\]
joins the identity to this map inside $SU_3$. It preserves the norm,
so it is a homotopy on the same local fibre. The induced cohomology
action is the identity. In the quadratic normal form, monodromy
is multiplication by $-1$ on four variables. Its action on the
degree-three vanishing cohomology is also $+1$.
\end{proof}

The algebraic global fibre $F_{\mathrm{alg}}=SL_3(\mathbb C)$ has
mixed Hodge groups $H^3=\mathbb Q(-2)$, $H^5=\mathbb Q(-3)$, and
$H^8=\mathbb Q(-5)$. To obtain these, use the algebraic
first-column map $SL_3\to\mathbb C^3\setminus0$.
Its fibre is $SL_2\times\mathbb C^2$, and Gaussian elimination gives
Zariski local trivializations. The base has $H^5=\mathbb Q(-3)$,
by the localization sequence at the origin. The fibre has
$H^3=\mathbb Q(-2)$. The same degree argument in the Leray spectral
sequence gives the stated Tate factors and their product.
The rational comparison with $F_0$ was proved above. A comparison
in a mixed-Hodge Hall theory must also specify the identification
of its local coefficient functor with this global-fibre convention.
The rational stalk theorem uses no such additional identification.

Conjugation acts trivially on these stalk cohomology groups.
At a scalar triple it acts on $SL_3$ by left multiplication through
the connected adjoint image of $GL_2$. At a nonzero rank-one point,
the stabilizer is connected and preserves the rank-one vanishing
cycle. This assertion does not compute equivariant extension data
between the stalk degrees or between the strata.

\begin{proposition}[A normal-fibre comparison]
\label{prop:c3-collision-fibre-map}
Along the rank-one path $M_t=\operatorname{diag}(t,0,0)$, $t>0$,
the block normal slice gives a comparison from a quadratic Milnor
fibre into the scalar Milnor fibre. On reduced rational cohomology,
restriction is an isomorphism in degree three and zero in degrees
five and eight.
\end{proposition}

\begin{proof}
Choose $\epsilon>0$ much smaller than $t$ and the local ball radii.
For $B\in SL_2$, set $\beta=(\epsilon/(2t))^{1/2}$.
The normal fibre embeds as
\[
 B\longmapsto\operatorname{diag}(t,\beta B),\qquad
 2\det\operatorname{diag}(t,\beta B)=\epsilon.
\]
Its compact $SU_2$ representative lies in the normal Milnor ball
when $\epsilon$ is sufficiently small, and in the larger scalar
Milnor ball when $t$ is sufficiently small. Dividing the matrix by
$\alpha=(\epsilon/2)^{1/3}$ identifies this embedding with
\[
 B\longmapsto g\operatorname{diag}(1,B),\qquad
 g=\operatorname{diag}(t/\alpha,\beta/\alpha,\beta/\alpha)\in SL_3.
\]
Left translation by $g$ is homotopic to the identity in $SL_3$.
The inclusion of the sufficiently large bounded fibre into $SL_3$
is a homotopy equivalence by the polar retraction just proved.
Hence the cohomology map is restriction along the block inclusion
$SU_2\hookrightarrow SU_3$. The spectral sequence in the preceding
proof identifies its degree-three restriction with an isomorphism.
The target has no cohomology in degrees five or eight.
\end{proof}

The stalk Euler characteristic is $1$ on both strata:
$1+1-1=1$ at a scalar triple. Thus this numerical invariant cannot
detect the two additional cohomology groups at the collision.
The normal-fibre restriction identifies the surviving degree-three
class. It does not specify all extension morphisms of $\mathcal P$.

\section{Flags at a Jordan collision}
\label{sec:c3-collision-flags}

Let $\mathfrak E$ classify a commuting triple together with a common
invariant line, and let $b:\mathfrak E\to\mathfrak M$ forget the line.
This morphism is proper: before taking the quotient, it is the
closed incidence subvariety in
$\operatorname{Crit}(f)\times\mathbb P^1$.
The second morphism $a:\mathfrak E\to\mathfrak M_1\times\mathfrak M_1$
records the line module and its quotient, where
$\mathfrak M_1=\mathbb C^3\times B\mathbb G_m$.

\begin{proposition}[The ramified incidence chart]
\label{prop:c3-collision-ramified}
Restrict to the open substack $\mathfrak U\subset\mathfrak M$ where
$A_0\ne0$. Every triple has unique expressions
$B_0=b_1A_0$ and $C_0=c_1A_0$.
If
\[
 A_0=\begin{pmatrix}h&e\\g&-h\end{pmatrix},\qquad
 \rho=h^2+eg=-\det A_0,
\]
then the flag morphism is the quotient by conjugation of the finite
flat double cover
\begin{equation}
 \widetilde U\longrightarrow U,
 \qquad \lambda^2=\rho,
 \quad U=\mathbb C^3\times
 (\mathfrak{sl}_2\setminus0)\times\mathbb C^2.
 \label{eq:c3-collision-double-cover}
\end{equation}
Both $U$ and $\widetilde U$ are smooth of dimension eight.
The involution is $\lambda\mapsto-\lambda$.
The scheme fibres over semisimple and nilpotent $A_0$ are respectively
two reduced points and $\operatorname{Spec}\mathbb C[\lambda]/(\lambda^2)$.
Over a scalar triple in the full stack, the fibre of $b$ is $\mathbb P^1$.
\end{proposition}

\begin{proof}
The centralizer of every nonzero traceless $2\times2$ matrix is
$\mathbb CI\oplus\mathbb CA_0$. This follows by solving its four
commutation equations on a chart with a nonzero matrix entry.
It gives the asserted unique coefficients, algebraically on those
charts. A common invariant line is therefore an eigenline of $A_0$.
If the eigenvalue is $\lambda$, the equation is $\lambda^2=\rho$.
Conversely, $A_0-\lambda I$ has rank one on this double cover.
Its kernel is a line bundle and gives the inverse morphism of
incidence schemes. A nonzero entry provides the local rank-one
minor needed for this assertion, including on nonreduced bases.

The cover algebra is free on $1,\lambda$. Its branch divisor on $U$
is $\rho=0$, whereas its ramification divisor on $\widetilde U$
is $\lambda=0$. The pullback of $\rho$ is $\lambda^2$.
The gradient of $\lambda^2-h^2-eg$ can vanish only
at $\lambda=h=e=g=0$, excluded by $A_0\ne0$. Thus the covering
space is smooth. The assertions about its fibres follow from the
quadratic equation.
For example, take $A_0=E$. On the line chart spanned by $(1,z)$,
invariance says
\[
 \det\bigl((1,z),E(1,z)\bigr)=-z^2=0.
\]
This gives the double point directly. At a scalar triple every
line is invariant, with no incidence equation, giving $\mathbb P^1$.
\end{proof}

\begin{proposition}[Oriented incidence transfer]
\label{prop:c3-collision-transfer}
On the smooth cover of Proposition~\ref{prop:c3-collision-ramified}, complex
orientation defines a trace and a pullback on rational cohomology,
\[
 T:H^*(\widetilde U,\mathbb Q)\longrightarrow H^*(U,\mathbb Q),
 \qquad r^*:H^*(U,\mathbb Q)\longrightarrow H^*(\widetilde U,\mathbb Q).
\]
They are conjugation-equivariant and descend to the quotient stacks.
They satisfy
\[
 Tr^*=2\operatorname{id},\qquad r^*T=1+\tau^*.
\]
At a branch point the trace on the single rational coefficient
stalk is multiplication by two. The algebraic trace pairing of
the finite-flat cover, in the basis $1,\lambda$, has matrix
\[
 \begin{pmatrix}2&0\\0&2\rho\end{pmatrix}.
\]
It is perfect precisely where $\rho$ is invertible.
\end{proposition}

\begin{proof}
For a proper map $r$ between smooth complex manifolds of the same
dimension, the adjunction $Rr_*r^!\mathbb Q\to\mathbb Q$ and their
complex orientations identify $r^!\mathbb Q$ with $\mathbb Q$.
This defines $T$. On the unramified locus it sums the sheet values.
Near the branch divisor, $d\rho\ne0$, so the map has analytic form
$(\lambda,z)\mapsto(\lambda^2,z)$.
Its local oriented degree is two. Its trace at the branch is
therefore twice the value on the unique preimage.

The involution preserves complex orientation. On the unramified
locus its norm sums the two sheets. At a branch point it is the
identity, so its norm is multiplication by two. These descriptions
give the two identities on constructible sheaves and hence on
cohomology. The orientation and adjunction are invariant under
conjugation. The same construction on Borel approximations therefore
gives the quotient-stack statement.

For the algebraic trace, multiplication by $1$ has trace two and
multiplication by $\lambda$ has matrix
$\left(\begin{smallmatrix}0&\rho\\1&0\end{smallmatrix}\right)$.
Thus $\operatorname{Tr}(\lambda)=0$ and
$\operatorname{Tr}(\lambda^2)=2\rho$. The pairing determinant is
$4\rho$. At a nilpotent point its radical is generated by $\lambda$.
\end{proof}

The incidence trace is a constructed extension of the separated
cover trace on this open chart. Its source is the flag space.
At the branch, the map $a$ records two equal point supports and
forgets a nontrivial extension. On stabilizers it sends
$sI+tN$ to $(s,s)$, with additive kernel. Thus $a$ is no longer
the ordered-support isomorphism used for separated points.
A critical Hall comparison must identify its coefficient morphism
with this oriented incidence transfer. The trace identities alone
do not construct that identification or a coproduct into the full
tensor product of charge-one critical complexes.

The spatial branch parameter $\rho=-\det A_0$ is a function on
representations. It differs from a discriminant of equivariant
Chern classes. The perfect pairing on two separated sheet fibres,
the finite-flat algebraic trace pairing, and a scalar Hall Hopf
pairing have different domains.

\section{The obstruction at the scalar stratum}
\label{sec:c3-collision-obstruction}

\begin{proposition}[Failure of a diagonal critical slice]
\label{prop:c3-collision-diagonal-failure}
Restricting $A_0$ to the diagonal line $\mathbb CH$ does not give
a critical slice at a scalar triple. The restricted potential
has critical points that are not critical in the full representation
space, and it omits nonzero nilpotent triples.
\end{proposition}

\begin{proof}
Write $A_0=xH$ and use arbitrary traceless $B_0,C_0$. Then
\[
 f|_{A_0=xH}=2x(b_Ec_F-b_Fc_E).
\]
This is a five-variable potential with two additional unused
diagonal coordinates. At $x=0$, $B_0=H$, $C_0=E$, all its partial
derivatives vanish. In the full space, however,
\[
 df(\dot A=F,\dot B=0,\dot C=0)
 =\operatorname{Tr}(F[H,E])=2.
\]
Scaling $B_0,C_0$ gives such examples arbitrarily near the scalar
triple. Conversely, $(A_0,B_0,C_0)=(E,0,0)$ is a full critical point
and cannot be conjugated to a diagonal $A_0$.
At a scalar triple the infinitesimal gauge map
$\xi\mapsto([\xi,A],[\xi,B],[\xi,C])$ is zero. It cannot remove
the two omitted off-diagonal directions of $A_0$.
\end{proof}

The full flag correspondence remains proper at a scalar triple,
but its positive-dimensional fibre prevents the finite-cover
construction from applying there. Proper base change for constant
coefficients gives the fibre cohomology
\[
\begin{array}{c|c}
\text{triple}&H^*(b^{-1}(x),\mathbb Q)\\\hline
\text{two distinct point supports}&\mathbb Q^2\text{ in degree }0\\
\text{nontrivial punctual extension}&\mathbb Q\text{ in degree }0\\
\text{scalar triple}&\mathbb Q\text{ in degrees }0,2.
\end{array}
\]
Constant coefficients forget the nonreduced structure of the
middle fibre. The target critical complex also changes at the
scalar stratum, by Theorem~\ref{thm:c3-scalar-milnor}.
It cannot be extended there as a rank-one constant complex in
degree $-4$: its stalk has additional groups in degrees $-2$ and $1$.
Their Euler contributions cancel, so a scalar count cannot repair
this discrepancy.

In the point-brane interpretation, a nilpotent tuple records a
nontrivial extension of two coincident branes. At the scalar tuple
all such extension directions become available, and the full
determinant replaces the transverse quadratic potential.
The Milnor calculation identifies the additional graded vanishing
cohomology. It does not identify these classes with primitives of
a Hall coproduct or with renormalized field-theoretic observables.

Extending the critical comparison through this stratum requires
the actual coefficient morphism for the two maps $a,b$, including
its extension data, Tate factors, and orientation transport.
A global coproduct must then have a specified target, support
condition, and localization. Compatibility with a continuous scalar
pairing is a further condition for forming a geometric Hall double.
The determinant calculation and the ramified incidence maps provide
the local data against which those constructions can be tested.
